\chapter{Literature Review}
\label{ch:litreview}
\label{lit:chapter}

This chapter places SmartLoad inside the existing technology and research
landscape. SmartLoad is adaptive load-management middleware: an NGINX data
plane sits in front of a backend pool, while an out-of-band decision plane
(anomaly detection, demand forecasting, and learned routing) and a control
plane work together over a Redis bus, with a deterministic classical scheduler
kept in place as a safety floor. To explain why we made each of these choices,
the chapter surveys the products and tools in the same space
(\cref{lit:sec:market}), reviews the state of the art in the research behind
each engine (\cref{lit:sec:research}), records the standards and best practices
the system follows and the areas where no standard yet exists
(\cref{lit:sec:standards}), looks at the data-protection duties that come with
collecting request telemetry (\cref{lit:sec:laws}), maps the design onto the
established codes of engineering ethics (\cref{lit:sec:ethics}), and closes by
pulling together what the survey teaches and the openings SmartLoad is built to
fill (\cref{lit:sec:conclusion}).

%======================================================================
\section{Market Survey}
\label{lit:sec:market}

The market that SmartLoad enters is mature in its individual layers but
fragmented across them. Three kinds of established product are relevant here:
the reverse proxies and load balancers that form the request-serving data
plane; the orchestrators and autoscalers that adjust capacity; and the
observability stack that supplies the telemetry on which any adaptive decision
must rest. SmartLoad does not try to replace any of these. It puts a subset of
them together and adds the decision layer that none of them ships. This section
walks through each kind, names the specific tools SmartLoad uses, and explains
the choices.

\subsection{Reverse proxies and load balancers}
\label{lit:sec:market-lb}

The main production load-balancing fabric is built on long-lived reverse
proxies that pick an upstream per request from a small set of well-understood
scheduling rules: round-robin, weighted round-robin, least-connections, and
consistent hashing. The three open-source implementations that define this
space are NGINX \cite{nginx}, HAProxy \cite{haproxy}, and Envoy \cite{envoy}.
All three publish the same family of scheduling rules together with passive and
active health checking.

\Cref{lit:tab:classical-lb} sets out those four rules, the decision each one
makes, and the condition under which it stops serving traffic well. The pattern
across the column is consistent: every rule decides from in-memory state alone,
so none of them can re-weight an upstream from external telemetry. In NGINX,
dynamic upstream weights need either a paid module, a configuration rewrite
followed by a reload, or an external operator process driving the change. The
data plane in all three projects is designed around in-memory static
configuration plus passive health probing, so telemetry-driven routing is out
of scope by construction. Cloud-managed equivalents (AWS Application Load
Balancer, GCP Cloud Load Balancing, Azure Front Door) expose the same
scheduling rules behind managed control planes, but they constrain weight
updates to coarse autoscaling-group events and offer no per-second decision
surface for an external controller to drive.

\begin{table}[ht]
\centering
\caption{Classical load-balancing scheduling rules as published by NGINX
\cite{nginx}, HAProxy \cite{haproxy}, and Envoy \cite{envoy}. Every rule
decides from in-memory state; none consumes an external telemetry feedback
loop.}
\label{lit:tab:classical-lb}
\small
\begin{tabularx}{\linewidth}{@{}lXX@{}}
\toprule
\headrow
\textbf{Rule} & \textbf{Decision} & \textbf{What it cannot do} \\
\midrule
Round-robin (RR) &
Hand requests to upstreams in fixed cyclic order. &
Assumes uniform capacity; a slow backend keeps taking $1/N$ of traffic until
passive health checks trip. \\
\addlinespace
Weighted RR (WRR) &
RR with a static per-upstream weight set at admit time. &
Weights are operator-set and rarely revisited, so the split drifts away from
real capacity over time. \\
\addlinespace
Least-connections (LC) &
Route to the upstream with the fewest in-flight connections. &
In-flight count is a noisy load proxy; it tracks capacity only when per-request
cost is uniform. \\
\addlinespace
IP / consistent hash &
Hash a request attribute (client IP, cookie) to a fixed slot. &
Sticky by design; cannot react to backend health inside the affinity window
without breaking session identity. \\
\bottomrule
\end{tabularx}
\end{table}

SmartLoad picks \textbf{NGINX} as its default data plane for three reasons.
First, it is everywhere. NGINX is among the most widely deployed reverse
proxies, which lowers the cost of putting SmartLoad in front of an existing
estate. Second, the hot-reload model fits. NGINX applies a rewritten upstream
block on \texttt{nginx -s reload} without dropping live connections, which is
exactly the actuation surface SmartLoad's control plane needs to nudge weights
on the order of seconds. Third, the upstream model is simple and explicit. The
upstream block is plain text that an out-of-band process can rewrite
deterministically and audit afterwards, which matters for the operator-override
property described in \cref{lit:sec:ethics}.

The data plane is deliberately \emph{pluggable} rather than NGINX-locked. If
the decision plane called \texttt{nginx -s reload} directly, the whole stack
would become NGINX-only, and the first HAProxy or Envoy user would force a
rewrite. So SmartLoad instead defines a \texttt{LoadBalancerAdapter} interface,
with one folder per adapter, each self-contained with its implementation,
documentation, and conformance tests. The NGINX adapter works in the prototype;
Envoy, HAProxy, and AWS Application Load Balancer adapters exist as scaffolded
stubs that a conformance suite proves interchangeable once implemented. The
data plane is therefore a replaceable component behind a stable contract, not a
hard dependency.

\slpitch{\textbf{Where SmartLoad sits.} SmartLoad does not replace NGINX. It
runs NGINX as its data plane and adds an out-of-band sidecar that consumes
Redis-published routing decisions, rewrites \texttt{upstream.conf}, and triggers
\texttt{nginx -s reload} on the order of seconds. The classical scheduler stays
in place as the deterministic safety floor and is never deprecated.}

\subsection{Service meshes}
\label{lit:sec:market-mesh}

Service meshes such as Istio \cite{istio} push routing, retries, mutual TLS,
and fine-grained traffic-shaping policy into a sidecar proxy attached to every
workload. This is powerful: a mesh can express canary splits, circuit-breaking,
and per-route policy uniformly across a fleet. The cost is operational weight.
A mesh adds a control plane, a per-pod sidecar, and a policy surface that has to
be learned, operated, and kept current, and for many teams the extra value does
not justify that overhead. SmartLoad targets the case where a team wants
adaptive, predictive load management in front of a modest backend pool without
adopting a full mesh. Where a mesh is already present, SmartLoad's
adapter-based data plane is meant to sit alongside it rather than compete with
the mesh's own routing layer.

\subsection{Orchestration and autoscaling}
\label{lit:sec:market-autoscale}

The orchestration layer that overlaps most directly with SmartLoad is the
Kubernetes Horizontal Pod Autoscaler (HPA) \cite{k8shpa}, together with
event-driven extensions such as KEDA \cite{keda} that scale on external signals
(queue depth, stream lag, custom metrics). These come from a well-documented
lineage of cluster-management systems, namely Borg, Omega, and Kubernetes
\cite{burns2016borgomega}, whose evolution shaped the reactive-scaling model in
use today. \Cref{lit:tab:autoscalers} compares the dominant production shapes.

The defining feature shared across the table is that the trigger is
\textbf{reactive}: a metric crosses a threshold, a stabilisation window prevents
flapping, and only then does the replica count move. The lag between load
arriving and a new replica being ready is bounded below by container start time
plus warm-up, commonly 30 to 90 seconds for typical web workloads. The same
order of magnitude appears in the large-scale fleet measurements that motivated
the Borg work \cite{verma2015borg}. Cloud autoscalers therefore react
\emph{after} load has already risen, so the spike is felt by users before
capacity arrives.

Predictive variants do exist. The autoscaling survey by Lorido-Botran et al.\
\cite{lorido2014autoscaling} catalogues them, and AWS Predictive Scaling is the
closest commercial analogue: it fits a forecaster on CloudWatch metrics and
pre-warms instances on the predicted shape. The forecaster is a black box, and,
more importantly, the forecasting and routing decisions stay decoupled, so the
prediction cannot inform the per-request scheduler.

\begin{table}[ht]
\centering
\caption{Production autoscaler shapes. The trigger is reactive in every case
except AWS predictive scaling, and the reactive lag is structurally bounded
below by container start plus warm-up (commonly 30 to 90 seconds).}
\label{lit:tab:autoscalers}
\small
\begin{tabularx}{\linewidth}{@{}lXlc@{}}
\toprule
\headrow
\textbf{System} & \textbf{Trigger} & \textbf{Scales} & \textbf{Posture} \\
\midrule
Kubernetes HPA \cite{k8shpa} &
CPU/memory or custom metric over a stabilisation window &
Replica count & Reactive \\
\addlinespace
KEDA \cite{keda} &
External event signals (queue depth, stream lag, custom) &
Replica count (HPA-backed) & Reactive \\
\addlinespace
K8s Cluster Autoscaler &
Pending unschedulable pods &
Node count & Reactive \\
\addlinespace
AWS Auto Scaling (target/step) &
CloudWatch metric breach &
EC2 instance count & Reactive \\
\addlinespace
AWS Predictive Scaling &
Built-in forecaster on historical metrics &
EC2 instance count & Proactive (decoupled) \\
\addlinespace
\textbf{SmartLoad autoscaler} &
Forecasted RPS, with reactive sliding window as fallback &
Backend pool count & Proactive + reactive floor \\
\bottomrule
\end{tabularx}
\end{table}

This reactive-versus-proactive split is well established in the literature.
Roy et al.\ \cite{roy2011autoscaling} show that provisioning ahead of demand
from a workload forecast measurably cuts SLA violations against purely reactive
control, while the Lorido-Botran survey \cite{lorido2014autoscaling} records
that reactive threshold scaling still dominates practice because it is simple
and dependable. SmartLoad's response is to make prediction first-class and to
couple it to actuation. A forecasting service publishes a short-horizon
request-rate forecast, and an autoscaler uses that forecast to grow the backend
pool \emph{ahead} of saturation, falling back to a reactive sliding-window rule
when the forecaster is silent or unhealthy. The contribution here is not a new
autoscaling algorithm. It is the proactive-by-default posture combined with a
reactive safety net.

\subsection{Observability stack}
\label{lit:sec:market-observability}

Any adaptive decision is only as good as the telemetry beneath it, so SmartLoad
uses a conventional, vendor-neutral observability stack rather than inventing
one. Request telemetry is emitted in OpenTelemetry's OTLP format
\cite{opentelemetry}, collected, and stored; per-process service health is
exposed in the Prometheus exposition format \cite{prometheus} and scraped
directly; operators view both surfaces in Grafana \cite{grafana}. Aggregated
time-series telemetry, namely per-backend latency, error rate, and request
count, lands in TimescaleDB \cite{timescaledb}, a PostgreSQL extension whose
hypertables suit the append-heavy, time-windowed query pattern that every
decision engine runs. The control plane coordinates over a Redis pub/sub bus
\cite{carlson2013redis}, chosen because the control signals are fire-and-forget
envelopes for which a lightweight broker is enough and a durable queue would be
over-engineered. The whole system is packaged and shipped with Docker
\cite{merkel2014docker} so an evaluator can stand the stack up locally and
reproduce results. Each of these is the de facto open standard in its niche, so
choosing them keeps SmartLoad portable and lets the new work concentrate in the
decision plane.

\subsection{Comparison}
\label{lit:sec:market-comparison}

\Cref{lit:tab:market} compares SmartLoad against the three closest alternatives
across the axes that matter for predictive load management. NGINX alone is the
deterministic data plane SmartLoad builds on, Kubernetes HPA is the dominant
reactive autoscaler, and a service mesh represents the heavyweight, policy-rich
routing fabric. SmartLoad is the only column that combines predictive scaling
and learned routing with a deterministic fallback and an operator override, and
it does so at a far lower operational weight than a mesh.

\begin{table}[ht]
\centering
\caption{SmartLoad compared with the closest alternatives across the axes
relevant to predictive, safety-bounded load management.}
\label{lit:tab:market}
\small
\begin{tabularx}{\linewidth}{@{}Y *{4}{>{\centering\arraybackslash}X}@{}}
\toprule
\headrow \textbf{Capability} & \textbf{SmartLoad} & \textbf{NGINX alone} & \textbf{K8s HPA} & \textbf{Service mesh} \\
\midrule
Predictive (forecast-driven) scaling & Yes & No & Reactive only & No \\
Learned (adaptive) routing            & Yes & No & No & Policy-driven, not learned \\
Deterministic fallback                & Yes (first-class) & Yes (only mode) & N/A & Partial \\
Operator override / kill switch       & Yes (per decision) & Config edit & Manual & Config edit \\
Operational weight                    & Low to moderate & Low & Moderate & High \\
\bottomrule
\end{tabularx}
\end{table}

%======================================================================
\section{Existing Research Efforts}
\label{lit:sec:research}

Where \cref{lit:sec:market} surveyed shipping products, this section surveys the
research behind each of SmartLoad's engines. The guiding research question
across the whole system is whether a load balancer can be made to adapt to
changing backend conditions, predicting load and learning routing preferences,
without giving up the deterministic guarantees that make classical load
balancing trustworthy. Each subsection below covers the slice of the literature
relevant to one engine and records what that literature establishes and where
its findings stop short of an integration-ready system.

\Cref{lit:tab:lineage} maps the foundations SmartLoad builds on before the
engine-by-engine discussion. It groups the literature by the engine it informs,
names the canonical reference in each line, and records what SmartLoad takes
from it. The remaining subsections expand each row in turn.

\begin{table}[ht]
\centering
\caption{Related-work lineage: the research foundations grouped by the
SmartLoad engine each informs, with the canonical reference and what SmartLoad
adopts from it.}
\label{lit:tab:lineage}
\small
\begin{tabularx}{\linewidth}{@{}llX@{}}
\toprule
\headrow
\textbf{Engine / concern} & \textbf{Canonical reference} & \textbf{What SmartLoad takes} \\
\midrule
Autoscaling (predictive) &
\cite{lorido2014autoscaling, roy2011autoscaling} &
Forecast-ahead provisioning cuts SLA breaches; reactive scaling stays the
dependable default. \\
\addlinespace
Time-series forecasting &
\cite{box2015timeseries, hyndman2021forecasting} &
ARIMA as the reference baseline; a simple model with uncertainty intervals and
a moving-average floor. \\
\addlinespace
Anomaly detection &
\cite{liu2008isolationforest, su2019omnianomaly} &
Unsupervised, label-free, small-artifact scoring; a labelled dataset used only
for an F1 gate. \\
\addlinespace
RL for systems &
\cite{sutton2018rl, schulman2017ppo, raffin2021sb3} &
Policy-gradient routing with a clipped objective and action masking as a hard
safety property. \\
\addlinespace
Autonomic computing &
\cite{kephart2003autonomic, ibm2005mapek} &
The MAPE-K loop (monitor, analyse, plan, execute over shared knowledge) as the
control-loop shape. \\
\addlinespace
Cluster traces &
\cite{verma2015borg, tirmazi2020borg, guo2019alibaba} &
Production-like workload for offline RL training and comparative benchmarks. \\
\bottomrule
\end{tabularx}
\end{table}

\subsection{Autoscaling and predictive capacity management}
\label{lit:sec:research-autoscale}

The autoscaling literature is broad and has been surveyed widely. The review by
Lorido-Botran et al.\ \cite{lorido2014autoscaling} sorts auto-scaling
techniques for cloud environments along the reactive--predictive and
threshold--model axes, and its main finding is that reactive, threshold-based
scaling dominates practice precisely because it is simple and dependable, while
predictive approaches promise lower tail latency at the cost of forecast risk.
The predictive branch is shown well by Roy et al.\ \cite{roy2011autoscaling},
who drive resource provisioning from a workload-forecasting model and a
look-ahead controller, and show that provisioning \emph{ahead} of demand
measurably cuts SLA violations compared with purely reactive control. The
research question SmartLoad takes from this line is how to get a trustworthy
short-horizon forecast cheaply, and how to degrade gracefully when the forecast
is wrong. That question motivates both the forecasting engine and its mandatory
reactive fallback.

\subsection{Time-series forecasting foundations}
\label{lit:sec:research-forecasting}

The forecasting engine rests on classical time-series theory. The Box--Jenkins
methodology for autoregressive integrated moving-average (ARIMA) modelling
\cite{box2015timeseries} gives the canonical framework for fitting an
autoregressive model to a differenced series and selecting its order from
autocorrelation structure; it is still the reference baseline against which
heavier forecasters are judged. The modern, practice-oriented synthesis by
Hyndman and Athanasopoulos \cite{hyndman2021forecasting} sets out the
principles SmartLoad's forecasting design follows: prefer a simple model with
clear uncertainty intervals, always keep a naive or moving-average baseline as
a floor, and evaluate on held-out data with a scale-free error metric. The key
finding the field offers is that, at short horizons, a well-specified ARIMA or
even a moving-average predictor is often competitive with much heavier sequence
models. That finding is why SmartLoad ships a moving-average baseline as a
first-class engine and an ARIMA bundle as the upgrade, rather than reaching
straight for a deep model.

\subsection{Anomaly detection}
\label{lit:sec:research-anomaly}

The anomaly-detection engine has to flag an unhealthy backend from its telemetry
without labelled training data, which points to unsupervised methods. Isolation
Forest \cite{liu2008isolationforest} is the canonical choice: it isolates
anomalies by random recursive partitioning, using the fact that anomalous
points sit in sparse regions of feature space and are therefore isolated in
fewer splits than normal points. Its training cost is linear, it needs no
labels, and it serialises to a small artifact. Those properties all fit a
routing-path engine that must fall back cleanly when its model is absent. For
the multivariate, contextual case, the line of work around the Server Machine
Dataset (SMD) and stochastic recurrent models \cite{su2019omnianomaly} shows
that sequence-aware detectors catch contextual anomalies that point-based
methods miss, and, importantly for evaluation, supplies a public dataset with
per-timestep ground-truth labels. SmartLoad uses that labelled dataset not for
training but to check detector quality against a fixed F1 gate. The field's
finding is a trade-off SmartLoad takes seriously. Point-based detectors are
cheap and explainable but blind to context, while sequence models are accurate
but expensive to train and serve, which is why the threshold rule is the
shipping baseline and the heavier detector is a guarded upgrade path.

\subsection{Reinforcement learning for routing and resource management}
\label{lit:sec:research-rl}

The learned-routing engine draws on the reinforcement-learning (RL) literature.
The foundational text by Sutton and Barto \cite{sutton2018rl} frames sequential
decision-making as a Markov Decision Process and develops the value- and
policy-based methods the field builds on. Among policy-gradient methods,
Proximal Policy Optimization (PPO) \cite{schulman2017ppo} is the common choice.
Its clipped surrogate objective turns off the gradient when the policy update
leaves a trust region, which makes it forgiving of mis-tuned learning rates
without an explicit constraint, and that robustness is what makes it practical
for a small team. SmartLoad uses the action-masking PPO variant from
Stable-Baselines3 \cite{raffin2021sb3} so that unhealthy backends are removed
from the action space before the policy ever samples, a hard safety property
rather than a learned preference.

SmartLoad applies reinforcement learning to routing as a learned policy that
sits behind proven deterministic baselines, with a safety fallback always in
place. SmartLoad frames the problem as an offline contextual bandit over
replayed traces \cite{guo2019alibaba} rather than an online Markov decision
process, which keeps training reproducible and bounds its cost. The learned
router runs in shadow mode by default, the classical scheduler stays the trusted
floor, and the anomaly constraints filter the action space before the policy is
ever consulted. This keeps the adaptive component bounded by guarantees the
operator can rely on while the policy is exercised against real traffic.

The honest evaluation finding, developed in the results chapter, is that on a
homogeneous-latency training distribution the PPO policy converges to
round-robin equivalence, since uniform sampling is the optimal policy when no
backend discriminates on latency. The policy is therefore kept in shadow mode
rather than promoted, and retraining on a heterogeneous-latency distribution is
the named next step. Richer state features, online updating, and a wider set of
replayed traces are the further research directions the project records.

\subsection{Autonomic computing and control loops}
\label{lit:sec:research-mapek}

SmartLoad's overall shape, which is to observe telemetry, decide, actuate, and
repeat, is an instance of the autonomic-computing vision set out by Kephart and
Chess \cite{kephart2003autonomic}, in which self-managing systems take low-level
control off the operator's hands. The concrete pattern is the MAPE-K loop
(Monitor, Analyse, Plan, Execute, over shared Knowledge) codified in IBM's
autonomic blueprint \cite{ibm2005mapek}. SmartLoad's mapping is direct: the
observability stack monitors, the decision engines analyse and plan, the
load-balancer sidecar and autoscaler execute, and the policy plus telemetry
store is the shared knowledge. Framing the architecture this way connects it to
a well-established body of work on safe, layered control and makes clear that
the new part is not the loop itself but the safety hierarchy imposed on its
plan stage.

\subsection{Cluster traces as data sources}
\label{lit:sec:research-traces}

Realistic evaluation of any load-management system needs realistic workload
data, and the field has settled on a small set of public production traces.
Google's Borg traces \cite{verma2015borg, tirmazi2020borg} and Alibaba's
cluster traces \cite{guo2019alibaba} document request patterns, resource usage,
and machine heterogeneity at fleet scale, and have become the standard data
sources for autoscaling and scheduling research. SmartLoad replays Alibaba
trace windows to drive both the offline RL training and the comparative
benchmarks, which grounds its claims in measured production-like load rather
than synthetic curves. The published trace studies also supply the empirical
basis for the tens-of-seconds provisioning-lag figure that motivates predictive
scaling in the first place.

%======================================================================
\section{Best Practices and Standards}
\label{lit:sec:standards}

SmartLoad adopts an established standard at every interface where one exists,
and is clear about the one place where no standard yet does. Treating these
standards as constraints rather than afterthoughts is what keeps the system
portable and interoperable with the surrounding observability and orchestration
ecosystem surveyed in \cref{lit:sec:market}.

The HTTP control surface, namely the operator API for reading state and posting
policy changes, is specified against OpenAPI 3.1 \cite{openapi}, so the
contract is machine-readable, client SDKs can be generated, and the interface
can be validated in continuous integration. Telemetry follows the OpenTelemetry
data model and is shipped over OTLP \cite{opentelemetry}, the emerging
cross-vendor standard for traces, metrics, and logs; this is what lets
SmartLoad's request telemetry flow into any conformant collector rather than a
proprietary agent. Per-process health is published in the Prometheus exposition
format \cite{prometheus}, the de facto standard for pull-based metrics, so any
Prometheus-compatible scraper can read service metrics without bespoke glue.

Beyond these wire formats, SmartLoad follows two widely adopted engineering
conventions. The first is semantic versioning, applied consistently to the HTTP
API, the policy schema, the Redis envelope, and the client SDK so that consumers
can reason about compatibility. The second is the twelve-factor application
methodology, whose discipline around configuration-in-environment, stateless
processes, and disposability maps cleanly onto a containerised,
horizontally-composed service set. For the eventual rollout of learned routing
into live traffic, the progressive-delivery patterns embodied by tools such as
Flagger \cite{flagger}, which provide canary analysis, automated promotion, and
automated rollback gated on metrics, are the reference model SmartLoad's shadow-
and hybrid-mode rollout follows.

There is, however, an open opportunity the project targets directly. No settled
standard yet covers learned routing with a safety fallback. The observability
and API layers are standardised, and progressive delivery standardises
\emph{how} a new behaviour is rolled out, but neither one says \emph{what
guarantees} an adaptive routing decision must satisfy before it is trusted, nor
how a deterministic fallback should be composed with a learned policy. SmartLoad
fills that space with its own contract, a strict routing decision hierarchy in
which anomaly constraints are applied first, learned recommendations second, and
the classical scheduler third, with an operator kill switch that short-circuits
all adaptation. This contract is a project convention, not an external standard.
Naming the opening is itself a contribution of the survey, since the safety
semantics of adaptive routing are an area where practice currently runs ahead of
standardisation.

%======================================================================
\section{Related Laws and Regulations}
\label{lit:sec:laws}

Any system that collects telemetry from live request traffic has to be checked
against data-protection law. The most influential regime is the European Union's
General Data Protection Regulation (GDPR) \cite{gdpr2016}, whose principles of
lawful basis, purpose limitation, data minimisation, and accountability now
shape similar regimes worldwide and are a reasonable benchmark even for a system
not deployed in the EU.

The relevant question is whether SmartLoad processes \emph{personal data}.
SmartLoad's decision plane works on \emph{operational metadata}: per-backend
request latencies, HTTP status codes, aggregate request counts, and internal
backend identifiers, grouped into time windows. It does not collect request
bodies, client identifiers, authentication tokens, or payload contents, and the
backend identifiers it records are internal pool labels rather than end-user
identities. On that basis the telemetry SmartLoad itself processes falls outside
the core definition of personal data, which materially reduces its regulatory
exposure.

That conclusion does not end the project's responsibilities. First, the boundary
is a design property that has to be maintained. If a future feature logged client
IP addresses or request headers, for example to support consistent-hash routing
on a client attribute, that data could become personal data, and the
data-minimisation and purpose-limitation principles would then apply directly.
So the design treats ``operational metadata only'' as an invariant to defend,
not a permanent guarantee. Second, a deployer who places SmartLoad in front of a
service that \emph{does} handle personal data inherits the controller's
obligations for the surrounding system, so SmartLoad's documentation should make
clear what it does and does not record so the deployer can complete their own
assessment. Finally, the project is developed under an academic accreditation
context whose criteria \cite{abet2025criteria} require explicit consideration of
the legal and societal impact of engineering work, which is the duty this
section discharges.

%======================================================================
\section{Related Code of Ethics}
\label{lit:sec:ethics}

SmartLoad's design decisions are not only technical; several of them are direct
expressions of professional engineering ethics. Two codes are authoritative
here: the ACM/IEEE-CS Software Engineering Code of Ethics and Professional
Practice \cite{gotterbarn1999ethics} and the ACM Code of Ethics and
Professional Conduct \cite{acm2018ethics}. Both place the public interest above
all other considerations, and SmartLoad's safety-first architecture can be read
as a concrete attempt to honour specific principles from them.

\paragraph{Public: reliability and safety.}
The first principle of the Software Engineering Code holds that engineers shall
act consistently with the public interest, which in an infrastructure component
means putting reliability first and avoiding harm to the systems that depend on
it. SmartLoad acts on this by refusing to let any adaptive decision remove the
deterministic safety floor. The classical scheduler is always present, the
decision plane sits out of the critical request path so that a failed engine
stops \emph{adaptation} rather than \emph{traffic}, and the routing hierarchy
guarantees that an unhealthy backend is excluded before any learned
recommendation is considered. The system is designed to fail safe, not only to
perform well when everything works.

\paragraph{Product: testing and quality.}
The Code's product principle requires that work meet the highest professional
standards, which for SmartLoad is enforced through pluggable engines that each
ship with a baseline requiring no model artifact, conformance test suites that
prove adapters interchangeable, and a discipline of unit tests and validation
gates before a change is accepted. Quality is treated as a property of the
artifact, established by tests, not a claim made in prose.

\paragraph{Judgment: transparency in reporting.}
The principle on professional judgment, backed by the ACM Code's commitment to
honesty and to avoiding misleading claims, is the reason SmartLoad reports its
evaluation results plainly and keeps the learned router in shadow mode while the
classical scheduler remains the trusted floor. An engineering report that
dressed up its adaptive component beyond what the measurements support would
break this principle directly. Reporting clearly, and letting the safety
hierarchy follow from the evidence, is how the project honours it.

\paragraph{Transparency and operator agency.}
Both codes value transparency and respect for the people affected by a system.
SmartLoad encodes this in two mechanisms. The first is an operator override, a
single kill-switch flag that immediately disables all adaptation and reverts to
deterministic routing. The second is an audit log that records every decision
the system actuates, surfaced through a transparent decision feed. Together the
audit log, the operator override, and that decision feed give the human
operator, not the model, final authority, and make every automated action
inspectable after the fact. This puts the Code's emphasis on honesty and
accountability into working machinery.

%======================================================================
\section{Conclusion}
\label{lit:sec:conclusion}

The survey establishes that the load-management space is strong in its
individual layers but leaves a clear gap between them. Classical load balancers
(\cref{lit:sec:market-lb}) provide deterministic, well-understood scheduling but
cannot react to backend health within seconds without external orchestration.
Cloud autoscalers (\cref{lit:sec:market-autoscale}) react to threshold breaches
but, outside narrow proprietary offerings, cannot pre-warm capacity against a
predicted spike, and even when they forecast they keep the prediction decoupled
from routing. The research on RL for routing (\cref{lit:sec:research-rl})
demonstrates feasibility but does not ship as integration-ready middleware, and
the commercial AIOps platforms surface anomalies and forecasts to a human while
deliberately declining to automate the data-plane response. No single existing
system closes the loop from detection to actuation while keeping that actuation
provably safe.

From this, SmartLoad's technical approach follows. It runs a proven classical
data plane (NGINX, behind a pluggable adapter) and adds an out-of-band decision
plane covering anomaly detection, forecasting, and learned routing, coordinated
over a Redis control bus and grounded in a conventional, standards-based
observability stack. The decision plane is composed with the data plane through
a strict safety hierarchy and an operator kill switch, so adaptiveness is only
ever exercised inside bounds the operator can revoke.

Three openings identified by the survey define the contribution. The first is to
predict ahead instead of reacting. SmartLoad drives scaling from a short-horizon
forecast and routing from current-state prediction, moving capacity before the
spike rather than after the alarm, with a reactive rule as the floor. The second
is to treat the deterministic fallback as a first-class design property. Every
engine ships with a dependency-free baseline and degrades to it cleanly, so the
system never enters a ``no decision available'' state and never loses its
deterministic guarantee. The third is the absence of a standard for safe learned
routing. The survey found mature standards for telemetry, APIs, and progressive
delivery, but none for the safety semantics of an adaptive routing decision, and
SmartLoad's routing hierarchy is offered as a concrete filling of that gap
pending an external standard. The chapters that follow develop each of these
into the system's design and evaluation.
